\documentclass[12pt,a4paper]{article}
\usepackage[utf8]{inputenc} % inutile avec XeLaTeX/LuaLaTeX
\usepackage[T1]{fontenc}
\usepackage{amsmath,amssymb,mathrsfs,tikz,times,pifont}
\usepackage{enumitem}
\usepackage{multicol}
\usepackage{lmodern}
\newcommand\circitem[1]{%
\tikz[baseline=(char.base)]{
\node[circle,draw=gray, fill=red!55,
minimum size=1.2em,inner sep=0] (char) {#1};}}
\newcommand\boxitem[1]{%
\tikz[baseline=(char.base)]{
\node[fill=cyan,
minimum size=1.2em,inner sep=0] (char) {#1};}}
\setlist[enumerate,1]{label=\protect\circitem{\arabic*}}
\setlist[enumerate,2]{label=\protect\boxitem{\alph*}}
\everymath{\displaystyle}
\usepackage[left=1cm,right=1cm,top=1cm,bottom=1.7cm]{geometry}
\usepackage[colorlinks=true, linkcolor=blue, urlcolor=blue, citecolor=blue]{hyperref}
\usepackage{array,multirow}
\usepackage[most]{tcolorbox}
\usepackage{varwidth}
\usepackage{float}
\tcbuselibrary{skins,hooks}
\usetikzlibrary{patterns}

\newtcolorbox{exa}[2][]{enhanced,breakable,before skip=2mm,after skip=5mm,
colback=yellow!20!white,colframe=black!20!blue,boxrule=0.5mm,
attach boxed title to top left ={xshift=0.6cm,yshift*=1mm-\tcboxedtitleheight},
fonttitle=\bfseries,
title={#2},#1,
boxed title style={frame code={
\path[fill=tcbcolback!30!black]
([yshift=-1mm,xshift=-1mm]frame.north west)
arc[start angle=0,end angle=180,radius=1mm]
([yshift=-1mm,xshift=1mm]frame.north east)
arc[start angle=180,end angle=0,radius=1mm];
\path[left color=tcbcolback!60!black,right color = tcbcolback!60!black,
middle color = tcbcolback!80!black]
([xshift=-2mm]frame.north west) -- ([xshift=2mm]frame.north east)
[rounded corners=1mm]-- ([xshift=1mm,yshift=-1mm]frame.north east)
-- (frame.south east) -- (frame.south west)
-- ([xshift=-1mm,yshift=-1mm]frame.north west)
[sharp corners]-- cycle;
},interior engine=empty,
},interior style={top color=yellow!5}}

\usepackage{fancyhdr}
\usepackage{eso-pic}
\usepackage{tkz-tab}
\AddToShipoutPicture{
    \AtTextCenter{%
        \makebox[0pt]{\rotatebox{80}{\textcolor[gray]{0.7}{\fontsize{5cm}{5cm}\selectfont PGB}}}
    }
}

\usepackage{verbatim}

\usepackage{color,soul}

\usepackage{amsmath}
\usepackage{amsfonts}
\usepackage{amssymb}
\usepackage{systeme}
\usepackage{tkz-tab}
\usepackage{tikz}
\usetikzlibrary{arrows}
\newcounter{exemple} % Compteur pour les questions

% Définir la commande pour afficher une question numérotée
\newcommand{\exemple}{%
  \refstepcounter{exemple}%
  \textbf{\textcolor{green}{Exemple \theexemple :}} \ignorespaces
}
%---------------------------------------
\definecolor{myorange}{rgb}{1.0, 0.8, 0.0}

% Définir un compteur pour les exercices d'application
\newcounter{exerciceapp}

% Définir la commande pour afficher un exercice d'application numéroté
\newcommand{\exerciceapp}{%
  \refstepcounter{exerciceapp}%
  \textbf{\textcolor{myorange}{Exercice d'application \theexerciceapp :}} \ignorespaces
}
%--------------------------------------
% Définir un compteur pour les exercices d'application
\newcounter{correction}

% Définir la commande pour afficher un correction exercice d'application numéroté
\newcommand{\correction}{%
  \refstepcounter{correction}%
  \textbf{\textcolor{myorange}{Correction \thecorrection :}} \ignorespaces
}
%--------------------------------------
% Commande pour ajouter du texte en arrière-plan
\usepackage{fancyhdr}
\usepackage{eso-pic}
%\usepackage{tkz-tab}
\AddToShipoutPicture{
    \AtTextCenter{%
        \makebox[0pt]{\rotatebox{80}{\textcolor[gray]{0.7}{\fontsize{5cm}{5cm}\selectfont PGB}}}
    }
}
%This command takes a colour as an optional argument; the default colour is black.
\usetikzlibrary{shapes.geometric,fit}
\newcommand{\myul}[2][black]{\setulcolor{#1}\ul{#2}\setulcolor{black}}
\newcommand\tab[1][1cm]{\hspace*{#1}}

\begin{document}
% En-tête personnalisée
\begin{center}
    \Large\textbf{\underline{\textcolor{red}{Fonction Exponentielle}}}\\[-0.1cm]
    \normalsize\textbf{Prof : M. BA} \hfill \textbf{Classe : TS2}\\[-0.1cm]
    \textbf{Année scolaire : 2025 -- 2026}
\end{center}

\subsection*{\underline{\textbf{\textcolor{red}{1. Définition}}}}

La fonction logarithme népérien \( \ln \) est continue et strictement croissante sur \( ]0;+\infty[ \). Elle réalise donc une bijection de \( ]0;+\infty[ \) sur \( \mathbb{R} \).

Elle admet alors une fonction réciproque, notée \( \exp \), appelée \textbf{fonction exponentielle}. Cette fonction est définie sur \( \mathbb{R} \) et à valeurs dans \( ]0;+\infty[ \).

On note :
\(
\forall x\in\mathbb{R},\qquad \exp(x)=e^x.
\)

Par définition, pour tout réel \(x\) et tout réel strictement positif \(y\),
\(
e^x=y
\quad\Longleftrightarrow\quad
x=\ln(y).
\)
\subsection*{\underline{\textbf{\textcolor{red}{2. Conséquences de la définition}}}}

\begin{enumerate}[label=(\alph*)]

\item \textbf{Ensemble de définition et image}

\(
D_{e^x}=\mathbb{R}
\qquad\text{et}\qquad
\mathrm{Im}(e^x)=]0;+\infty[.
\)

En particulier,

\(
\forall x\in\mathbb{R},\qquad e^x>0.
\)

\item \textbf{Relations avec le logarithme}

Pour tout réel \(x\) et tout réel \(a>0\),

\(
e^{\ln(a)}=a
\)

et

\(
\ln(e^x)=x.
\)

\item \textbf{Valeurs remarquables}

\(
e^0=1
\qquad;\qquad
e^1=e\simeq2,718.
\)

\item \textbf{Sens de variation}

La fonction exponentielle est continue et strictement croissante sur \( \mathbb{R} \).

Par conséquent,

\(
a<b
\Longrightarrow
e^a<e^b.
\)

\end{enumerate}

\subsection*{\underline{\textbf{\textcolor{red}{3. Propriétés algébriques}}}}

\underline{\textbf{\textcolor{red}{Propriété fondamentale}}}

Pour tous réels \(a\) et \(b\),

\(
e^{a+b}=e^a\times e^b.
\)

Cette propriété permet d'établir les égalités suivantes.

\underline{\textbf{\textcolor{red}{Conséquences}}}

Pour tous réels \(a\), \(b\) et \(r\),

\begin{itemize}
\item \(
e^0=1,
\)

\item \(
e^{-a}=\frac1{e^a},
\)

\item \(
e^{a-b}=\frac{e^a}{e^b},
\)

\item \(
e^{ra}=(e^a)^r.
\)
\end{itemize}

Comme la fonction exponentielle est strictement croissante,

\(
e^a=e^b
\Longleftrightarrow
a=b,
\)

et

\(
e^a<e^b
\Longleftrightarrow
a<b.
\)

\subsection*{\underline{\textbf{\textcolor{red}{4. Limites}}}}

\textbf{Aux bornes de l'ensemble de définition}

\begin{enumerate}
\item \(
\lim\limits_{x\to-\infty}e^x=0
\)

\item \(
\lim\limits_{x\to+\infty}e^x=+\infty
\)
\end{enumerate}

\textbf{Limites usuelles}
\begin{enumerate}
\item \(
\lim\limits_{x\to+\infty}\frac{e^x}{x}=+\infty
\)

\item \(
\lim\limits_{x\to-\infty}xe^x=0
\)

\item \(
\lim\limits_{x\to0}\frac{e^x-1}{x}=1.
\)
\end{enumerate}
\underline{\textbf{\textcolor{red}{Preuve de quelques limites}}}\\

\underline{\textbf{\textcolor{red}{Exercice d'application}}}

Calculer les limites suivantes :
\begin{enumerate}
\item \(
\lim\limits_{x\to+\infty}\frac{3e^x-2}{5e^x+3}
\)
\item \(
\lim\limits_{x\to-\infty}\frac{\ln(1+e^x)}{e^x}
\)
\item \(
\lim\limits_{x\to+\infty}(x-e^x)
\)
\item \(
\lim\limits_{x\to+\infty}(x-e^x)
\)
\item \(
\lim\limits_{x\to+\infty}\frac{\sin(2x)}{1-e^x}.
\)
\end{enumerate}

\subsection*{\underline{\textbf{\textcolor{red}{5. Limites des fonctions composées avec l'exponentielle}}}}

\underline{\textbf{\textcolor{red}{Propriété}}}

Soit \(u\) une fonction définie sur un intervalle \(I\).

\begin{itemize}
    \item Si \(\displaystyle\lim\limits_{x\to a}u(x)=+\infty\), alors
    \(
    \lim\limits_{x\to a}e^{u(x)}=+\infty.
    \)

    \item Si \(\displaystyle\lim\limits_{x\to a}u(x)=-\infty\), alors
    \(
    \lim\limits_{x\to a}e^{u(x)}=0.
    \)

    \item Si \(\displaystyle\lim\limits_{x\to a}u(x)=\ell\in\mathbb{R}\), alors
    \(
    \lim\limits_{x\to a}e^{u(x)}=e^{\ell}.
    \)
\end{itemize}

\underline{\textbf{\textcolor{red}{Méthode}}}

Pour calculer une limite contenant une exponentielle composée :

\begin{enumerate}
    \item On détermine d'abord la limite de l'exposant \(u(x)\).
    \item On utilise ensuite les limites usuelles de la fonction exponentielle.
\end{enumerate}

\underline{\textbf{\textcolor{red}{Exemple}}}

Calculer la limite suivante :
\(
\lim\limits_{x\to+\infty}e^{-3x+5}.
\)

\underline{\textbf{\textcolor{red}{Solution}}}

On commence par calculer la limite de l'exposant :
\(
\lim\limits_{x\to+\infty}(-3x+5)=-\infty.
\)

Or,
\(
\lim\limits_{t\to-\infty}e^t=0.
\)

En posant \(t=-3x+5\), on obtient
\(
\boxed{\lim\limits_{x\to+\infty}e^{-3x+5}=0.}
\)

\underline{\textbf{\textcolor{red}{Exercices d'application}}}

Calculer les limites suivantes :

\begin{enumerate}
\item \(
\lim\limits_{x\to+\infty}e^{2x-1}
\)

\item
\(
\lim\limits_{x\to-\infty}e^{-5x+4}
\)

\item
\(
\lim\limits_{x\to+\infty}e^{\frac{1-x}{2}}
\)

\item
\(
\lim\limits_{x\to0}e^{\sin x}
\)

\item
\(
\lim\limits_{x\to+\infty}e^{\frac{3x+1}{x-2}}.
\)
\end{enumerate}

Ces cinq exemples couvrent les trois situations classiques : un exposant qui tend vers \(+\infty\), vers \(-\infty\) ou vers une limite réelle finie.

\subsection*{\underline{\textbf{\textcolor{red}{6. Dérivée d'une fonction exponentielle composée}}}}

\underline{\textbf{\textcolor{red}{Propriété}}}

Soit \(u\) une fonction dérivable sur un intervalle \(I\).

Alors la fonction
\(
x\longmapsto e^{u(x)}
\)

est dérivable sur \(I\) et
\(
\left(e^{u(x)}\right)'=u'(x)e^{u(x)}.
\)

\underline{\textbf{\textcolor{red}{Notation}}}

La composée \(\exp\circ u\) est généralement notée \(e^u\).

Sa dérivée est donc
\(
(e^u)'=u'e^u.
\)

\underline{\textbf{\textcolor{red}{Exemple}}}\\
Déterminer les limites suivantes:\\
\begin{itemize}
\item[] $\bullet$ La fonction $x \longmapsto e^{-x^{2}+x}$ est dérivable sur $\mathbb{R}$ et sa dérivée est la fonction \( \boxed{f'(x)=(-2x+1)e^{-x^2+x}} \) \\
\item[] $\bullet$ La fonction $x \longmapsto e^{\cos x}$ est dérivable sur $\mathbb{R}$ et sa dérivée est  la fonction \( \boxed{g'(x)=-\sin(x)e^{\cos x}} \)\\
\item[] $\bullet$ La fonction $x \longmapsto e^{\frac{1}{x}}$ est dérivable sur $\mathbb{R}^{*}$ et sa dérivée est  la fonction \( \boxed{h'(x)=-\frac1{x^2}e^{\frac1x}} \)\\
\end{itemize}

\subsection*{\underline{\textbf{\textcolor{red}{8. Croissance comparée de \(\ln x\), \(x^\alpha\) et \(e^x\)}}}}

\underline{\textbf{\textcolor{red}{a. Introduction}}}

Lorsqu'une variable \(x\) tend vers \(+\infty\), certaines fonctions
croissent beaucoup plus rapidement que d'autres.

On dit qu'une fonction \(f\) a une croissance plus rapide qu'une fonction
\(g\) lorsque le quotient \(\frac{f(x)}{g(x)}\) tend vers \(+\infty\).

La comparaison des fonctions usuelles conduit aux résultats suivants :
\(
\ln x \ll x^\alpha \ll e^x
\)

où \(\alpha\) est un réel strictement positif.

Cela signifie que :
\begin{itemize}
    \item les fonctions puissances \(x^\alpha\) dominent le logarithme ;
    \item la fonction exponentielle domine toutes les fonctions puissances.
\end{itemize}

\underline{\textbf{\textcolor{red}{b. Croissance comparée de \(\ln x\) et \(x^\alpha\)}}}

Soit \(\alpha>0\). On a :
\(
\boxed{
\lim\limits_{x\to+\infty}\frac{\ln x}{x^\alpha}=0
}
\)

Autrement dit :
\(
\ln x \text{ croît moins rapidement que } x^\alpha.
\)

On a également :
\(
\boxed{
\lim\limits_{x\to0^+}x^\alpha\ln x=0
}
\)

\underline{\textbf{\textcolor{red}{c. Croissance comparée de \(x^\alpha\) et \(e^x\)}}}

Soit \(\alpha>0\). Alors :
\(
\boxed{
\lim\limits_{x\to+\infty}\frac{x^\alpha}{e^x}=0
}
\)

Ainsi :
\(
x^\alpha \text{ croît moins rapidement que } e^x.
\)

\underline{\textbf{\textcolor{red}{d. Résumé des croissances comparées}}}

Lorsque \(x\to+\infty\),\(
\boxed{
\ln x \ll x^\alpha \ll e^x
}
\)
\begin{itemize}
\item
\(
\lim\limits_{x\to+\infty}\frac{\ln x}{x^\alpha}=0,
\)
\item
\(
\lim\limits_{x\to+\infty}\frac{x^\alpha}{e^x}=0.
\)

\item
\(
\lim\limits_{x\to+\infty}\frac{\ln x}{e^x}=0.
\)
\end{itemize}
\subsection*{\underline{\textbf{\textcolor{red}{9. Équations, systèmes et inéquations faisant intervenir l'exponentielle}}}}

\subsubsection*{\underline{\textbf{\textcolor{red}{a°) Équations exponentielles}}}}

\underline{\textbf{\textcolor{red}{Méthodes}}}

Pour résoudre une équation contenant une exponentielle, on utilise généralement l'une des propriétés suivantes :

\begin{itemize}
    \item \(e^x>0\) pour tout \(x\in\mathbb{R}\).
    \item \(e^a=e^b \Longleftrightarrow a=b\).
    \item \(e^x=a\) avec \(a>0\Longleftrightarrow x=\ln(a)\).
\end{itemize}

\underline{\textbf{\textcolor{red}{Exemple}}}

Résoudre dans \(\mathbb{R}\) les équations suivantes :

\begin{enumerate}[label=\alph*)]
    \item \(e^x=-1\)
    \item \(e^{x+1}=3\)
    \item \(e^{x^2}=e^{x+2}\)
    \item \((e^x-2)(e^{-x}+1)=0\)
\end{enumerate}

\subsubsection*{\underline{\textbf{\textcolor{red}{ b°) Systèmes d'équations }}}}

\underline{\textbf{\textcolor{red}{Méthode}}}

Dans un système contenant des exponentielles, on pose souvent

\(
X=e^x
\quad\text{et}\quad
Y=e^y,
\)

avec

\(
X>0,\qquad Y>0.
\)

Le système devient alors un système algébrique classique.

\underline{\textbf{\textcolor{red}{Exemples}}}

Résoudre les systèmes suivants :

\(
\begin{cases}
4e^x-3e^y=9\\
2e^x+e^y=7
\end{cases}
\)

\(
\begin{cases}
e^xe^y=10\\
e^{x-y}=\dfrac25
\end{cases}
\)

\(
\begin{cases}
e^{2x}-7e^{y+1}=-10\\
x-y=1
\end{cases}
\)

\subsubsection*{\underline{\textbf{\textcolor{red}{ c°) Inéquations exponentielles }}}}

\underline{\textbf{\textcolor{red}{Méthodes}}}

Pour résoudre une inéquation exponentielle, on utilise le fait que la fonction exponentielle est strictement croissante :

\(
e^a<e^b
\Longleftrightarrow
a<b,
\)

\(
e^a\le e^b
\Longleftrightarrow
a\le b.
\)

Lorsque l'inéquation est polynomiale en \(e^x\), on effectue le changement de variable

\(
X=e^x
\qquad(X>0).
\)

\underline{\textbf{\textcolor{red}{Exemples}}}

Résoudre dans \(\mathbb{R}\) les inéquations suivantes :

\begin{enumerate}[label=\alph*)]
    \item \(e^{-x}\ge2\)
    \item \(e^{x^2-3}\le e^{2x}\)
    \item \(2e^{2x}-5e^x+2>0\)
\end{enumerate}
\section*{\underline{\textbf{\textcolor{red}{10.Etude le fonction exp}}}}
Soit $f(x)=exp(x)$ le domaine \\
Le Domaine $D_{f}$ \\
$D_{f}=\mathbb{R}$\\
$\bigotimes$ Limites aux bornes de $D_{f}$\\
\underline{En $-\infty $}\\
$\lim_{x \to -\infty} e^{x}=0$\\
\underline{En $+\infty$}\\
$\lim_{{x \to +\infty}} e^{x}=+\infty$\\
$\bigotimes$ La dérivée de f\\
$f'(x)=e^{x}$\\
$\forall x \in \mathbb{R}$, $f'(x)>0$, donc f est croissante sur $]0;+\infty[$\\
$\bigotimes$ Tableau de variation\\
%Tableau de Variation
\definecolor{cqcqcq}{rgb}{0.7529411764705882,0.7529411764705882,0.7529411764705882}
\begin{tikzpicture}[line cap=round,line join=round,>=triangle 45,x=1cm,y=1cm]
\draw [color=cqcqcq,, xstep=1cm,ystep=1cm] (-7,-10) grid (-22,17);
\clip(-22,-5) rectangle (12,10);
\draw [line width=2pt] (-23,8)-- (-7,8); %première ligne A(-22,8)---B(-7,8)
\draw [line width=2pt] (-22,6)-- (-7,6); %deuxième ligne
\draw [line width=2pt] (-22,4)-- (-7,4); %troisième ligne
\draw [line width=2pt] (-22,-2)-- (-7,-2);%dernière ligne
\draw [line width=2pt] (-22,-2)-- (-22,8); %première colonne
\draw [line width=2pt] (-19,8)-- (-19,-2); %deuxième colone
\draw [line width=2pt] (-7,8)-- (-7,-2); %troisième colonne
\draw (-21,1.5) node[anchor=north west] {$f(x)$};
\draw (-21,5.5) node[anchor=north west] {$f'(x)$};
\draw (-21,7) node[anchor=north west] {$x$};
\draw (-19,7) node[anchor=north west] {$-\infty$};
\draw (-8,7) node[anchor=north west] {$+\infty$};
%Asymptote verticale
%\draw [line width=2pt] (-18.7,6)-- (-18.7,-2);
%¨\draw [line width=2pt] (-18.79,6)-- (-18.79,-2);
%signe de la dérivé
\draw (-13.5,5.3) node[anchor=north west] {$+$};
%zéro de exp
\draw [line width=2pt] (-13,6)-- (-13,1);
\draw (-13.2,6.5) node[anchor=north west] {$0$};
\draw (-13.2,1.2) node[anchor=north west] {$1$};
\draw [->,line width=2pt] (-18,-1) -- (-8,3);
\draw (-18.5,-1) node[anchor=north west] {$0$};
\draw (-8,3.5) node[anchor=north west] {$+\infty$};
\end{tikzpicture}

\begin{tikzpicture}[scale=1.5]
  \draw[->] (-5,0) -- (2,0) node[right] {$x$};
  \draw[->] (0,-0.5) -- (0,2) node[above] {$y$};
  \draw[scale=1,domain=-5:1.5,smooth,variable=\x,blue] plot ({\x},{exp(\x)});
  %\node[right,blue] at (1.5,1.5) {$f(x) = e^x$};
  \draw[scale=1,domain=-1.5:1.5,smooth,variable=\x,blue] plot ({\x},{\x+1});
  
%  \draw[scale=1,domain=-5:1.5,smooth,variable=\x,blue] plot ({\x},{\ln(\x)});
  %\node[right,blue] at (1.5,1.5) {$f(x) = e^x$};
%  \draw[scale=1,domain=-1.5:1.5,smooth,variable=\x,blue] plot ({\x},{\x-1});
\end{tikzpicture}

$C_{f}$ est au-dessous de sa tangente en J; donc $\forall x \in \mathbb{R}, e^{x}>x+1$
\section*{\underline{\textbf{\textcolor{red}{11.Branche infinie de ln}}}}
On a  $\lim_{{x \to +\infty}} e^{x}=+\infty$ et $\lim_{{x \to +\infty}} \frac{e^{x}}{x}=+\infty$\\
Nous avons ainsi une branche parabolique de direction (Oy) au voisinage de +$\infty$.\\
Car $\lim_{{x \to +\infty}} \frac{ln(x)}{x}=0$\\
\section*{\underline{\textbf{\textcolor{red}{12.Application}}}}
\end{document}
